% GENERATED FILE. Edit canonical lessons, then run npm run generate:books.
% canonical-source-hash: fnv1a64:16c9d40b16c2f977
% canonical-lessons: HI-C483-taqat, HI-C483-cust, HI-C483-dam, HI-C483-bal, HI-C483-thakana, HI-R483-review-476-483

\chapter{Strength and Breath}
\label{ch:hi-strength-and-breath}
\hlchaptermodality{483}

\begin{chapteropening}
\textbf{By the end of this chapter:} \emph{I can say how strong or out of breath someone is.}

Complete the last lesson of chapter 483: I can say how strong or out of breath someone is.
\end{chapteropening}

\section[\texorpdfstring{\hi{ताक़त} (tāqat)}{ताक़त (tāqat)}]{\hi{ताक़त} (tāqat) — strength; power}

\label{lesson:HI-C483-taqat}



\begin{quote}
Before the new one: say the Hindi for fair-skinned, then the Hindi for lean.
\end{quote}

\subsection*{You'll want to know: \hi{ताक़त}}
\textbf{\hi{ताक़त}} — \emph{tāqat} — \textquotedblleft{}strength; power\textquotedblright{}.

\begin{quote}
\hi{दूध} \hi{पीने} \hi{से} \hi{बच्चों} \hi{को} \hi{ताक़त} \hi{मिलती} \hi{है।}
\end{quote}

\begin{quote}
\emph{dūdh pīne se baccoṁ ko tāqat miltī hai.}
\end{quote}

\begin{quote}
Drinking milk gives children strength.
\end{quote}

Feminine.

\begin{cousinweb}
\textbf{\hi{ताक़त}} (\emph{tāqat}) is usually traced to Arabic \emph{ṭāqa}, \textquotedblleft{}ability, strength\textquotedblright{}, which reached Hindi through Persian.
\end{cousinweb}

\begin{grammarlens}[title={What you've built}]
One more word for describing what you lived and felt.
\end{grammarlens}

\subsection*{Your turn}
\begin{itemize}
\raggedright
  \item \emph{Say it:} \emph{tāqat}
  \item \emph{Say it:} \emph{tāqat}, once more
  \item \emph{Say it:} \emph{dūdh pīne se baccoṁ ko tāqat miltī hai.} — Drinking milk gives children strength.
  \item \emph{Recall:} say the Hindi for fair-skinned, then the Hindi for lean, then say \emph{tāqat} again
\end{itemize}

\begin{tcolorbox}[breakable,colback=teal!4,colframe=teal!35!black,title={Before you move on}]
What does \hi{ताक़त} mean? (\textquotedblleft{}Strength; power\textquotedblright{}.) Say it once more.
\end{tcolorbox}
\section[\texorpdfstring{\hi{चुस्त} (cust)}{चुस्त (cust)}]{\hi{चुस्त} (cust) — spry; fit}

\label{lesson:HI-C483-cust}



\begin{quote}
Before the new one: say the Hindi for lean, then the Hindi for strength.
\end{quote}

\subsection*{You'll want to know: \hi{चुस्त}}
\textbf{\hi{चुस्त}} — \emph{cust} — \textquotedblleft{}spry; fit\textquotedblright{}.

\begin{quote}
\hi{अस्सी} \hi{की} \hi{उम्र} \hi{में} \hi{भी} \hi{नानी} \hi{बहुत} \hi{चुस्त} \hi{हैं।}
\end{quote}

\begin{quote}
\emph{assī kī umr meṁ bhī nānī bahut cust haiṁ.}
\end{quote}

\begin{quote}
Even at eighty, my grandmother is very spry.
\end{quote}

\begin{cousinweb}
\textbf{\hi{चुस्त}} (\emph{cust}) is usually traced to Persian \emph{chust}, \textquotedblleft{}quick, active, tight\textquotedblright{}. Clothes can be \emph{cust} too: close-fitting.
\end{cousinweb}

\begin{grammarlens}[title={What you've built}]
One more word for describing what you lived and felt.
\end{grammarlens}

\subsection*{Your turn}
\begin{itemize}
\raggedright
  \item \emph{Say it:} \emph{cust}
  \item \emph{Say it:} \emph{cust}, once more
  \item \emph{Say it:} \emph{assī kī umr meṁ bhī nānī bahut cust haiṁ.} — Even at eighty, my grandmother is very spry.
  \item \emph{Recall:} say the Hindi for lean, then the Hindi for strength, then say \emph{cust} again
\end{itemize}

\begin{tcolorbox}[breakable,colback=teal!4,colframe=teal!35!black,title={Before you move on}]
What does \hi{चुस्त} mean? (\textquotedblleft{}Spry; fit\textquotedblright{}.) Say it once more.
\end{tcolorbox}
\section[\texorpdfstring{\hi{दम} (dam)}{दम (dam)}]{\hi{दम} (dam) — stamina; breath}

\label{lesson:HI-C483-dam}



\begin{quote}
Before the new one: say the Hindi for strength, then the Hindi for spry.
\end{quote}

\subsection*{You'll want to know: \hi{दम}}
\textbf{\hi{दम}} — \emph{dam} — \textquotedblleft{}stamina; breath\textquotedblright{}.

\begin{quote}
\hi{दादा} \hi{में} \hi{अभी} \hi{भी} \hi{बहुत} \hi{दम} \hi{है।}
\end{quote}

\begin{quote}
\emph{dādā meṁ abhī bhī bahut dam hai.}
\end{quote}

\begin{quote}
Grandfather still has plenty of stamina.
\end{quote}

Masculine.

\begin{cousinweb}
\textbf{\hi{दम}} (\emph{dam}) is usually traced to Persian \emph{dam}, \textquotedblleft{}breath\textquotedblright{}, and so the strength to keep going. \emph{dam lenā} is to catch one's breath.
\end{cousinweb}

\begin{grammarlens}[title={What you've built}]
One more word for describing what you lived and felt.
\end{grammarlens}

\subsection*{Your turn}
\begin{itemize}
\raggedright
  \item \emph{Say it:} \emph{dam}
  \item \emph{Say it:} \emph{dam}, once more
  \item \emph{Say it:} \emph{dādā meṁ abhī bhī bahut dam hai.} — Grandfather still has plenty of stamina.
  \item \emph{Recall:} say the Hindi for strength, then the Hindi for spry, then say \emph{dam} again
\end{itemize}

\begin{tcolorbox}[breakable,colback=teal!4,colframe=teal!35!black,title={Before you move on}]
What does \hi{दम} mean? (\textquotedblleft{}Stamina; breath\textquotedblright{}.) Say it once more.
\end{tcolorbox}
\section[\texorpdfstring{\hi{बल} (bal)}{बल (bal)}]{\hi{बल} (bal) — force; might}

\label{lesson:HI-C483-bal}



\begin{quote}
Before the new one: say the Hindi for spry, then the Hindi for stamina.
\end{quote}

\subsection*{You'll want to know: \hi{बल}}
\textbf{\hi{बल}} — \emph{bal} — \textquotedblleft{}force; might\textquotedblright{}.

\begin{quote}
\hi{अरुण} \hi{ने} \hi{पूरे} \hi{बल} \hi{से} \hi{दरवाज़ा} \hi{धकेला।}
\end{quote}

\begin{quote}
\emph{Aruṇ ne pūre bal se darvāzā dhakelā.}
\end{quote}

\begin{quote}
Arun pushed the door with all his strength.
\end{quote}

Masculine.

\begin{cousinweb}
\textbf{\hi{बल}} (\emph{bal}) is inherited from Sanskrit \emph{bála}, \textquotedblleft{}strength, power\textquotedblright{}: CDIAL lists Hindi \emph{bal}, \textquotedblleft{}strength, ability\textquotedblright{}, under it (CDIAL 9161). \emph{dublā} is built on its opposite, \emph{durbala}.
\end{cousinweb}

\begin{grammarlens}[title={What you've built}]
One more word for describing what you lived and felt.
\end{grammarlens}

\subsection*{Your turn}
\begin{itemize}
\raggedright
  \item \emph{Say it:} \emph{bal}
  \item \emph{Say it:} \emph{bal}, once more
  \item \emph{Say it:} \emph{Aruṇ ne pūre bal se darvāzā dhakelā.} — Arun pushed the door with all his strength.
  \item \emph{Recall:} say the Hindi for spry, then the Hindi for stamina, then say \emph{bal} again
\end{itemize}

\begin{tcolorbox}[breakable,colback=teal!4,colframe=teal!35!black,title={Before you move on}]
What does \hi{बल} mean? (\textquotedblleft{}Force; might\textquotedblright{}.) Say it once more.
\end{tcolorbox}
\section[\texorpdfstring{\hi{थकाना} (thakānā)}{थकाना (thakānā)}]{\hi{थकाना} (thakānā) — to tire someone out}

\label{lesson:HI-C483-thakana}



\begin{quote}
Before the new one: say the Hindi for stamina, then the Hindi for force.
\end{quote}

\subsection*{You'll want to know: \hi{थकाना}}
\textbf{\hi{थकाना}} — \emph{thakānā} — \textquotedblleft{}to tire someone out\textquotedblright{}.

\begin{quote}
\hi{लंबे} \hi{सफ़र} \hi{ने} \hi{बच्चों} \hi{को} \hi{थका} \hi{दिया।}
\end{quote}

\begin{quote}
\emph{lambe safar ne baccoṁ ko thakā diyā.}
\end{quote}

\begin{quote}
The long journey tired the children out.
\end{quote}

Past, with \hi{ने} (he or she …): \hi{उसने} \hi{बच्चों} \hi{को} \hi{थकाया}, \emph{usne baccoṁ ko thakāyā}.

\begin{cousinweb}
\textbf{\hi{थकाना}} (\emph{thakānā}) is the causative partner of \emph{thaknā}, \textquotedblleft{}to become weary\textquotedblright{}, which CDIAL lists under \emph{*sthakk}, a form it reconstructs (CDIAL 13737).
\end{cousinweb}

\begin{grammarlens}[title={What you've built}]
One more word for describing what you lived and felt.
\end{grammarlens}

\subsection*{Your turn}
\begin{itemize}
\raggedright
  \item \emph{Say it:} \emph{thakānā}
  \item \emph{Say it:} \emph{thakānā}, once more
  \item \emph{Say it:} \emph{lambe safar ne baccoṁ ko thakā diyā.} — The long journey tired the children out.
  \item \emph{Recall:} say the Hindi for stamina, then the Hindi for force, then say \emph{thakānā} again
\end{itemize}

\begin{tcolorbox}[breakable,colback=teal!4,colframe=teal!35!black,title={Before you move on}]
What does \hi{थकाना} mean? (\textquotedblleft{}To tire someone out\textquotedblright{}.) Say it once more.
\end{tcolorbox}
\section[(dialogue)]{Looking back: from The Face to Strength and Breath}

\label{lesson:HI-R483-review-476-483}



\begin{quote}
Say the Hindi for a face and an eyebrow.
\end{quote}

\subsection*{Your turn}
Say these aloud:

\begin{itemize}
\raggedright
  \item a face, an eyebrow, a moustache, a beard, a jaw
  \item the palm of the hand, a fist, an arm, a leg, a heel
  \item a lung, the brain, a vein, a rib, a kidney
  \item the mouth, a gum (of the teeth), to suck, to chew, to spit
  \item a taste, a touch, a smell, an aroma, to lick
  \item to lie down, to draw in, to limp, to stand on tiptoe, to jump
  \item height, appearance, dusky, fair-skinned, lean
  \item strength, spry, stamina, force, to tire someone out
\end{itemize}

\begin{tcolorbox}[breakable,colback=teal!4,colframe=teal!35!black,title={Before you move on}]
A face? (\textbf{\hi{चेहरा}}, \emph{cehrā}.) An eyebrow? (\textbf{\hi{भौंह}}, \emph{bhauṁh}.) A moustache? (\textbf{\hi{मूँछ}}, \emph{mū̃ch}.) A beard? (\textbf{\hi{दाढ़ी}}, \emph{dāṛhī}.) A jaw? (\textbf{\hi{जबड़ा}}, \emph{jabṛā}.) The palm of the hand? (\textbf{\hi{हथेली}}, \emph{hathelī}.) A fist? (\textbf{\hi{मुट्ठी}}, \emph{muṭṭhī}.) An arm? (\textbf{\hi{बाँह}}, \emph{bā̃h}.) A leg? (\textbf{\hi{टाँग}}, \emph{ṭā̃g}.) A heel? (\textbf{\hi{एड़ी}}, \emph{eṛī}.) A lung? (\textbf{\hi{फेफड़ा}}, \emph{phephṛā}.) The brain? (\textbf{\hi{दिमाग़}}, \emph{dimāġ}.) A vein? (\textbf{\hi{नस}}, \emph{nas}.) A rib? (\textbf{\hi{पसली}}, \emph{paslī}.) A kidney? (\textbf{\hi{गुर्दा}}, \emph{gurdā}.) The mouth? (\textbf{\hi{मुँह}}, \emph{mũh}.) A gum (of the teeth)? (\textbf{\hi{मसूड़ा}}, \emph{masūṛā}.) To suck? (\textbf{\hi{चूसना}}, \emph{cūsnā}.) To chew? (\textbf{\hi{चबाना}}, \emph{cabānā}.) To spit? (\textbf{\hi{थूकना}}, \emph{thūknā}.) A taste? (\textbf{\hi{स्वाद}}, \emph{svād}.) A touch? (\textbf{\hi{स्पर्श}}, \emph{sparś}.) A smell? (\textbf{\hi{गंध}}, \emph{gandh}.) An aroma? (\textbf{\hi{महक}}, \emph{mahak}.) To lick? (\textbf{\hi{चाटना}}, \emph{cāṭnā}.) To lie down? (\textbf{\hi{लेटना}}, \emph{leṭnā}.) To draw in? (\textbf{\hi{सिकोड़ना}}, \emph{sikoṛnā}.) To limp? (\textbf{\hi{लँगड़ाना}}, \emph{lãgṛānā}.) To stand on tiptoe? (\textbf{\hi{उचकना}}, \emph{ucaknā}.) To jump? (\textbf{\hi{कूदना}}, \emph{kūdnā}.) Height? (\textbf{\hi{क़द}}, \emph{qad}.) Appearance? (\textbf{\hi{हुलिया}}, \emph{huliyā}.) Dusky? (\textbf{\hi{साँवला}}, \emph{sā̃vlā}.) Fair-skinned? (\textbf{\hi{गोरा}}, \emph{gorā}.) Lean? (\textbf{\hi{दुबला}}, \emph{dublā}.) Strength? (\textbf{\hi{ताक़त}}, \emph{tāqat}.) Spry? (\textbf{\hi{चुस्त}}, \emph{cust}.) Stamina? (\textbf{\hi{दम}}, \emph{dam}.) Force? (\textbf{\hi{बल}}, \emph{bal}.) To tire someone out? (\textbf{\hi{थकाना}}, \emph{thakānā}.)
\end{tcolorbox}
